\subsection{Esquemas regulatorios}

Esta subsección especifica los cinco diseños regulatorios considerados en el modelo. El esquema \(A_0\) constituye el punto de referencia institucional, mientras que \(A\), \(A'\), \(A''\) y \(B\) introducen modificaciones sucesivas en cuatro dimensiones: la composición del ICG, el universo utilizado para verificar el avance financiero, la regla de activación de las consecuencias pecuniarias y la forma de verificación aplicable a las intervenciones de \(R^{F}\). La tecnología de ejecución, la producción de resultados, el tratamiento del ahorro y las bases definidas en la Subsección~4.1 son comunes a todos los diseños.

La transición \(A_0\to A\) modifica simultáneamente varias dimensiones y se utiliza como \textit{benchmark} institucional. En cambio, la secuencia \(A\to A'\to A''\to B\) introduce cambios más acotados: \(A\to A'\) elimina el umbral agregado de activación; \(A'\to A''\) retira el avance financiero focalizado del ICG y sustituye su verificación por indicadores físicos o funcionales; y \(A''\to B\) concentra el ICG en los resultados del servicio. De acuerdo con la estructura de la Subsección~4.1, \(ICG^m\), \(\chi^m\) y \(M^m\) se evalúan para los niveles de cumplimiento correspondientes a cada esquema.

\subsubsection{Definición de los esquemas}

De acuerdo con la regla de promedio aritmético adoptada para el ICG, cada indicador individual incluido en un esquema recibe el mismo peso. Como \(C_S\) y \(C_G\) son, respectivamente, los promedios de los indicadores de \(J^{S}\) y \(J^{G}\), sus ponderaciones externas dependen del número de indicadores de cada grupo que permanecen dentro del índice. Sea \(N^m\) el número total de indicadores individuales incorporados en el ICG del esquema \(m\). Dado que \(n_S\equiv|J^{S}|\) y \(n_G\equiv|J^{G}|\), definidos en la Subsección~4.1, se tiene

\begin{equation}
\label{eq:numero_componentes_icg}
\begin{aligned}
N^{A_0}
&=
n_S+n_G+2,
&
N^{A}
=
N^{A'}
&=
n_S+n_G+1,
\\
N^{A''}
&=
n_S+n_G,
&
N^{B}
&=
n_S.
\end{aligned}
\end{equation}

Los dos componentes adicionales de \(A_0\) son el avance financiero general \(f_R\) y la relación de trabajo \(c_{RT}\); \(A\) y \(A'\) mantienen únicamente el avance financiero focalizado \(f_F\) además de los indicadores de \(J^{S}\) y \(J^{G}\). En \(A''\) permanecen solo estos dos grupos de indicadores y, en \(B\), únicamente los resultados del servicio.

La representación agregada es equivalente al promedio de los indicadores individuales. Para \(H\in\{S,G\}\),

\begin{equation}
\label{eq:equivalencia_promedio_bloques}
\frac{1}{N^m}
\sum_{j\in J^{H}}c_j
=
\frac{n_H}{N^m}C_H.
\end{equation}

Por tanto, las medidas agregadas de desempeño quedan definidas por

\begin{subequations}
\label{eq:icg_esquemas}
\begin{align}
ICG^{A_0} &= \frac{n_S}{N^{A_0}}C_S+\frac{n_G}{N^{A_0}}C_G+\frac{1}{N^{A_0}}f_R+\frac{1}{N^{A_0}}c_{RT},
\label{eq:icg_a0_esquemas}\\
ICG^{A} &= \frac{n_S}{N^{A}}C_S+\frac{n_G}{N^{A}}C_G+\frac{1}{N^{A}}f_F,
\label{eq:icg_a}\\
ICG^{A'} &= ICG^{A},
\label{eq:icg_ap}\\
ICG^{A''} &= \frac{n_S}{N^{A''}}C_S+\frac{n_G}{N^{A''}}C_G,
\label{eq:icg_app}\\
ICG^{B} &= C_S.
\label{eq:icg_b}
\end{align}
\end{subequations}

Para el canal reputacional definido en la Subsección~4.1, se distingue esta ponderación regulatoria de la composición del desempeño percibido \(P^m\). En las comparaciones del modelo se supone que los bloques que integran la señal pública corresponden a los componentes de desempeño incorporados en cada diseño, de modo que

\begin{equation}
\label{eq:bloques_reputacionales_esquemas}
\begin{aligned}
\mathrm{Sig}^{A_0}
&=
\{S,G,F,RT\},
&
\mathrm{Sig}^{A}
=
\mathrm{Sig}^{A'}
&=
\{S,G,F\},
\\
\mathrm{Sig}^{A''}
&=
\{S,G\},
&
\mathrm{Sig}^{B}
&=
\{S\}.
\end{aligned}
\end{equation}

Las ponderaciones reputacionales \(\alpha_H^m\) se obtienen, por tanto, a partir de las saliencias \(s_H\) definidas en la Subsección~4.1 y no de los cocientes \(n_H/N^m\). Estos últimos continúan siendo exclusivamente las ponderaciones regulatorias del ICG.

La reasignación de ponderaciones es, por tanto, una consecuencia mecánica de la salida de indicadores del ICG. En particular,

\begin{equation}
\label{eq:reasignacion_pesos_icg}
\frac{n_S}{N^{A}}
>
\frac{n_S}{N^{A_0}},
\qquad
\frac{n_S}{N^{A''}}
>
\frac{n_S}{N^{A'}},
\qquad
\frac{n_G}{N^{A''}}
>
\frac{n_G}{N^{A'}}.
\end{equation}

siempre que \(n_S>0\) y \(n_G>0\). Así, \(A_0\to A\) incrementa el peso regulatorio relativo de los resultados del servicio dentro del ICG; \(A'\to A''\) incrementa simultáneamente los pesos regulatorios relativos de \(C_S\) y \(C_G\); y \(A''\to B\) concentra el ICG en \(C_S\). La estática comparada de la Subsección~4.3 distinguirá estas ponderaciones regulatorias de las ponderaciones reputacionales \(\alpha_H^m\).

La verificación financiera se aplica únicamente en \(A_0\), \(A\) y \(A'\). De acuerdo con la notación de la Subsección~4.1,

\begin{equation}
\label{eq:verificacion_financiera_esquemas}
\begin{aligned}
f^{A_0}&=f_R,
&
\mathrm{Fin}^{A_0}&=R,
\\
f^{A}=f^{A'}&=f_F,
&
\mathrm{Fin}^{A}=\mathrm{Fin}^{A'}&=R^{F}.
\end{aligned}
\end{equation}

En \(A''\) y \(B\), las intervenciones de \(R^{F}\) se verifican mediante indicadores específicos de \(J^{F}\) sobre su cumplimiento físico o funcional. Para identificar compactamente la presencia de una base financiera se mantiene

\begin{equation}
\label{eq:indicador_base_financiera}
\eta_F^m
=
\begin{cases}
1,
&
m\in\{A_0,A,A'\},
\\[4pt]
0,
&
m\in\{A'',B\}.
\end{cases}
\end{equation}

Las reglas de activación son

\begin{subequations}
\label{eq:activacion_esquemas}
\begin{align}
\chi^{A_0}
&=
\mathbf{1}\{ICG^{A_0}<\tau\},
\label{eq:chi_a0}
\\
\chi^{A}
&=
\mathbf{1}\{ICG^{A}<\tau\},
\label{eq:chi_a}
\\
\chi^{A'}
&=
\chi^{A''}
=
\chi^{B}
=
1.
\label{eq:chi_sin_umbral}
\end{align}
\end{subequations}

En particular, \(A_0\) y \(A\) condicionan la consecuencia pecuniaria al umbral agregado del ICG, mientras que desde \(A'\) los incumplimientos configurados dejan de depender de esa condición adicional.

Combinando las bases definidas en la Subsección~4.1 con las reglas anteriores, las multas por esquema son

\begin{subequations}
\label{eq:multas_esquemas}
\begin{align}
M^{A_0}
&=
\chi^{A_0}\,
\mathrm{WACC}
\left(
B_S+B_G+B_{RT}+B_{\mathrm{fin}}^{A_0}
\right),
\label{eq:multa_base_a0}
\\
M^{A}
&=
\chi^{A}\,
\mathrm{WACC}
\left(
B_S+B_G+B_{\mathrm{fin}}^{A}
\right),
\label{eq:multa_a}
\\
M^{A'}
&=
\mathrm{WACC}
\left(
B_S+B_G+B_{\mathrm{fin}}^{A'}
\right),
\label{eq:multa_ap}
\\
M^{A''}
&=
\mathrm{WACC}
\left(
B_S+B_G+B_F
\right),
\label{eq:multa_app}
\\
M^{B}
&=
\mathrm{WACC}
\left(
B_S+B_G+B_F
\right).
\label{eq:multa_b}
\end{align}
\end{subequations}

En \(A_0\), una misma intervención de \(R^{S}\) o \(R^{G}\) puede incidir en la base física correspondiente a su propia meta, en la relación de trabajo y, además, en el avance financiero general. La transición \(A_0\to A\) elimina estas dos últimas fuentes de superposición para dichas intervenciones. Por su parte, la salida de \(f_F\) en \(A''\) no elimina la exigibilidad de \(R^{F}\), sino que sustituye la referencia financiera por los indicadores específicos de \(J^{F}\). Del mismo modo, la salida de \(C_G\) en \(B\) lo retira del ICG y, bajo la composición reputacional de \eqref{eq:bloques_reputacionales_esquemas}, concentra \(P^B\) en \(C_S\), pero no modifica la exigibilidad de las metas de \(J^{G}\).

\subsubsection{Exposición sancionadora por esquema}

Las diferencias anteriores se trasladan a la exposición sancionadora \(\mu_r^m\) definida en \eqref{eq:exposicion_sancionadora_r}. Esta exposición no constituye una multa adicional: representa cuánto aumenta la multa total cuando la intervención \(r\) pasa del estado de culminación al de no culminación, manteniendo constantes las demás condiciones relevantes.

Considérese una región regulatoria en la que la configuración de las demás obligaciones y la regla de activación \(\chi^m\) permanecen constantes entre ambos estados. Sea \(\delta_r^m\geq0\) el número de bases físicas activas en las que la no culminación de \(r\) incrementa la base sancionadora en \(\bar I_r\). Su valor es local y depende del esquema, de las obligaciones incumplidas y de las bases activas en el estado considerado.

Para evitar ambigüedad, los subíndices \(N\) y \(C\) denotan, respectivamente, el estado en que la inversión \(r\) está \emph{no culminada} (\(a_r<1\)) y el estado en que está \emph{culminada} (\(a_r=1\)), manteniendo constantes las demás decisiones y obligaciones. Así, \(B_{\mathrm{fin},N}^m\) y \(B_{\mathrm{fin},C}^m\) son los valores de la base financiera en ambos estados, mientras que \(M_{r,N}^m\) y \(M_{r,C}^m\) son las multas totales correspondientes.

Dentro de esa región, la exposición puede descomponerse como

\begin{subequations}
\label{eq:descomposicion_exposicion_esquemas}
\begin{align}
\mu_{r,\mathrm{fis}}^m &= \chi^m\,\mathrm{WACC}\,\bar I_r\,\delta_r^m,
\label{eq:exposicion_fisica_esquemas}\\
\mu_{r,\mathrm{fin}}^m &= \eta_F^m\,\chi^m\,\mathrm{WACC}\left(B_{\mathrm{fin},N}^m-B_{\mathrm{fin},C}^m\right),
\label{eq:exposicion_financiera_esquemas}\\
\mu_r^m &= \mu_{r,\mathrm{fis}}^m+\mu_{r,\mathrm{fin}}^m.
\label{eq:exposicion_total_esquemas}
\end{align}
\end{subequations}

El componente financiero es nulo cuando \(\eta_F^m=0\) o cuando \(r\notin\mathrm{Fin}^m\). La multiplicidad \(\delta_r^m\) recoge, en cambio, la posible superposición de bases físicas. En \(A_0\), por ejemplo, una intervención de \(R^{S}\) o \(R^{G}\) puede participar simultáneamente en la base de su propia meta y en \(B_{RT}\); desde \(A\), esta última fuente de superposición desaparece. Para \(R^{F}\), \(B_F\) en \(A''\) y \(B\) sustituye el mecanismo financiero utilizado en \(A\) y \(A'\).

La Tabla~\ref{tab:multiplicidad_bases_fisicas} muestra la multiplicidad física máxima que puede alcanzar una inversión en cada esquema. Estos valores son cotas superiores: el valor efectivo de \(\delta_r^m\) es local y depende de que las obligaciones correspondientes se encuentren configuradas como incumplidas en la región considerada.

\begin{table}[t]
\centering
\caption{Multiplicidad potencial de bases físicas por esquema}
\label{tab:multiplicidad_bases_fisicas}
\scriptsize
\renewcommand{\arraystretch}{1.25}
\setlength{\tabcolsep}{3pt}

\begin{tabularx}{\columnwidth}{
@{}
>{\centering\arraybackslash}p{0.13\columnwidth}
>{\centering\arraybackslash}p{0.24\columnwidth}
>{\centering\arraybackslash}p{0.20\columnwidth}
>{\raggedright\arraybackslash}X
@{}
}
\hline
\textbf{Esquema}
&
\textbf{\(r\in R^{S}\cup R^{G}\)}
&
\textbf{\(r\in R^{F}\)}
&
\textbf{Bases físicas potenciales}
\\
\hline

\(A_0\)
&
\(\delta_r^{A_0}\leq2\)
&
\(\delta_r^{A_0}\leq1\)
&
Meta propia y \(RT\) para \(R^{S}\cup R^{G}\); \(RT\) para \(R^{F}\).
\\[3pt]

\(A,\,A'\)
&
\(\delta_r^m\leq1\)
&
\(\delta_r^m=0\)
&
Meta propia para \(R^{S}\cup R^{G}\); \(R^{F}\) se verifica financieramente.
\\[3pt]

\(A'',\,B\)
&
\(\delta_r^m\leq1\)
&
\(\delta_r^m\leq1\)
&
Meta propia para \(R^{S}\cup R^{G}\) e indicador específico de \(J^{F}\) para \(R^{F}\).
\\

\hline
\end{tabularx}

\vspace{1mm}

\begin{minipage}{0.96\columnwidth}
\raggedright
\footnotesize
\textit{Nota:} Las cotas consideran únicamente bases físicas. La exposición total puede incluir además el componente financiero de \eqref{eq:exposicion_financiera_esquemas}. Cuando una obligación no se encuentra configurada como incumplida, su contribución a \(\delta_r^m\) es nula.
\end{minipage}

\end{table}

La descomposición muestra que los esquemas modifican la exposición sancionadora mediante tres canales: la multiplicidad de bases físicas \(\delta_r^m\), la presencia de una base financiera \(\eta_F^m\) y la regla de activación \(\chi^m\). Por ello, no se impone un ordenamiento general de \(\mu_r^m\): la comparación depende de la intervención, de las obligaciones configuradas y del estado regulatorio considerado.

En particular, \(A_0\to A\) reduce las fuentes potenciales de superposición para \(R^{S}\) y \(R^{G}\), al retirar estas intervenciones del avance financiero general y eliminar la relación de trabajo de la estructura sancionadora. Ello puede reducir \(\mu_r^m\), pero no permite inferir por sí solo un ordenamiento del esfuerzo, pues simultáneamente cambian la composición y las ponderaciones regulatorias del ICG, así como la composición y legibilidad de la señal reputacional \(P^m\), además de otros canales continuos de incentivo. Esta comparación se desarrolla en la Subsección~4.3.

La descomposición de \eqref{eq:descomposicion_exposicion_esquemas} requiere que \(\chi^m\) permanezca constante entre los estados de culminación y no culminación. En \(A_0\) y \(A\), la culminación de una intervención puede modificar el ICG y hacer que se cruce el umbral \(\tau\). Cuando ello ocurre, \(M_{r,N}^m-M_{r,C}^m\) incorpora además un cambio de régimen de activación y no puede representarse únicamente mediante las contribuciones física y financiera anteriores. Este efecto pertenece al margen discreto de activación definido en la Subsección~4.1 y se considera separadamente en las comparaciones de la Subsección~4.3.

\subsubsection{Resumen de los esquemas}

La Tabla~\ref{tab:resumen_esquemas_regulatorios} resume las diferencias institucionales que sustentan la estática comparada posterior. La columna \(N^m\) muestra el número de indicadores individuales que integran el ICG y determina únicamente sus ponderaciones regulatorias de los bloques agregados mediante \eqref{eq:equivalencia_promedio_bloques}; las ponderaciones reputacionales se rigen separadamente por \eqref{eq:ponderaciones_reputacionales}.

\begin{table*}[t]
\centering
\caption{Composición del ICG y consecuencias pecuniarias por esquema}
\label{tab:resumen_esquemas_regulatorios}
\scriptsize
\renewcommand{\arraystretch}{1.35}
\setlength{\tabcolsep}{2.0pt}

\begin{tabularx}{\textwidth}{
@{}
>{\centering\arraybackslash}p{0.045\textwidth}
>{\centering\arraybackslash}p{0.105\textwidth}
>{\centering\arraybackslash}p{0.225\textwidth}
>{\centering\arraybackslash}p{0.105\textwidth}
>{\centering\arraybackslash}p{0.145\textwidth}
>{\centering\arraybackslash}p{0.12\textwidth}
>{\centering\arraybackslash}X
@{}
}
\hline

\textbf{Esquema}
&
\textbf{\(N^m\)}
&
\textbf{\(ICG^m\)}
&
\textbf{Verificación financiera}
&
\textbf{Obligaciones con consecuencia}
&
\textbf{Activación}
&
\textbf{Consecuencia pecuniaria \(M^m\)}
\\

\hline

\(A_0\)
&
\(\displaystyle n_S+n_G+2\)
&
\(\displaystyle
\frac{n_S}{N^{A_0}}C_S
+
\frac{n_G}{N^{A_0}}C_G
+
\frac{f_R+c_{RT}}{N^{A_0}}
\)
&
\(R\)
&
\(J^{S},\,J^{G},\,RT,\,f_R\)
&
\(\displaystyle
\mathbf{1}\{ICG^{A_0}<\tau\}
\)
&
\(\displaystyle
\chi^{A_0}\mathrm{WACC}
\left(
B_S+B_G+B_{RT}+B_{\mathrm{fin}}^{A_0}
\right)
\)
\\[8pt]

\(A\)
&
\(\displaystyle n_S+n_G+1\)
&
\(\displaystyle
\frac{n_S}{N^{A}}C_S
+
\frac{n_G}{N^{A}}C_G
+
\frac{f_F}{N^{A}}
\)
&
\(R^{F}\)
&
\(J^{S},\,J^{G},\,f_F\)
&
\(\displaystyle
\mathbf{1}\{ICG^{A}<\tau\}
\)
&
\(\displaystyle
\chi^{A}\mathrm{WACC}
\left(
B_S+B_G+B_{\mathrm{fin}}^{A}
\right)
\)
\\[8pt]

\(A'\)
&
\(\displaystyle n_S+n_G+1\)
&
\(\displaystyle
ICG^{A}
\)
&
\(R^{F}\)
&
\(J^{S},\,J^{G},\,f_F\)
&
\(1\)
&
\(\displaystyle
\mathrm{WACC}
\left(
B_S+B_G+B_{\mathrm{fin}}^{A'}
\right)
\)
\\[8pt]

\(A''\)
&
\(\displaystyle n_S+n_G\)
&
\(\displaystyle
\frac{n_S}{N^{A''}}C_S
+
\frac{n_G}{N^{A''}}C_G
\)
&
---
&
\(J^{S},\,J^{G},\,J^{F}\)
&
\(1\)
&
\(\displaystyle
\mathrm{WACC}
\left(
B_S+B_G+B_F
\right)
\)
\\[8pt]

\(B\)
&
\(\displaystyle n_S\)
&
\(\displaystyle
C_S
\)
&
---
&
\(J^{S},\,J^{G},\,J^{F}\)
&
\(1\)
&
\(\displaystyle
\mathrm{WACC}
\left(
B_S+B_G+B_F
\right)
\)
\\

\hline
\end{tabularx}

\vspace{1mm}

\begin{minipage}{0.98\textwidth}
\raggedright
\footnotesize
\textit{Nota:} \(N^m\) corresponde al número de indicadores individuales incluidos en el ICG de cada esquema. Los coeficientes de \(C_S\) y \(C_G\) resultan de agregar, respectivamente, \(n_S\) y \(n_G\) indicadores con ponderación individual \(1/N^m\); estas ponderaciones son regulatorias y no deben confundirse con las ponderaciones reputacionales \(\alpha_H^m\). En \(A''\) y \(B\), la ausencia de verificación financiera indica que \(R^{F}\) se evalúa mediante indicadores específicos de \(J^{F}\) sobre su cumplimiento físico o funcional. La transición \(A_0\to A\) modifica simultáneamente la composición del ICG, el universo financiero y la estructura de consecuencias; las transiciones posteriores aíslan sucesivamente la activación por umbral, la verificación financiera y la concentración del ICG en los resultados del servicio.
\end{minipage}

\end{table*}
